\section{命题逻辑}

\subsection{逻辑变量和联结词}

逻辑变量是表示真/假的变量，取值为 $0$ 或 $1$。

其中，一个变量 $p$ 称为\textbf{逻辑变元}，而 $0, 1$ 称为\textbf{逻辑常元}。

联结词是逻辑变量之间的运算。下面是常用的联结词：

\subsubsection{逻辑非}

对于逻辑变量 $p$，其逻辑非记做 $\lnot p$，其真值表如下：

\begin{table}[htbp]
    \centering
    \begin{tabular}{|c|c|}
        \hline
        $p$ & $\lnot p$ \\
        \hline
        $0$ & $1$ \\
        \hline
        $1$ & $0$ \\
        \hline
    \end{tabular}
    \caption{逻辑非的真值表}
\end{table}

\subsubsection{合取}

对于逻辑变量 $p, q$，合取运算记做 $p \land q$，其真值表如下：

\begin{table}[htbp]
    \centering
    \begin{tabular}{|c|c|c|}
        \hline
        \diagbox{$p$}{$p \land q$}{$q$} & $0$ & $1$ \\
        \hline
        $0$ & $0$ & $0$ \\
        \hline
        $1$ & $0$ & $1$ \\
        \hline
    \end{tabular}
    \caption{合取的真值表}
\end{table}

\subsubsection{析取}

对于逻辑变量 $p, q$，析取运算记做 $p \lor q$，其真值表如下：

\begin{table}[htbp]
    \centering
    \begin{tabular}{|c|c|c|}
        \hline
        \diagbox{$p$}{$p \lor q$}{$q$} & $0$ & $1$ \\
        \hline
        $0$ & $0$ & $1$ \\
        \hline
        $1$ & $1$ & $1$ \\
        \hline
    \end{tabular}
    \caption{析取的真值表}
\end{table}

\subsubsection{异或}

对于逻辑变量 $p, q$，异或运算记做 $p \oplus q$，其真值表如下：

\begin{table}[htbp]
    \centering
    \begin{tabular}{|c|c|c|}
        \hline
        \diagbox{$p$}{$p \oplus q$}{$q$} & $0$ & $1$ \\
        \hline
        $0$ & $0$ & $1$ \\
        \hline
        $1$ & $1$ & $0$ \\
        \hline
    \end{tabular}
    \caption{异或的真值表}
\end{table}

\subsubsection{蕴涵}

对于逻辑变量 $p, q$，蕴涵运算记做 $p \to q$，其真值表如下：

\begin{table}[htbp]
    \centering
    \begin{tabular}{|c|c|c|}
        \hline
        \diagbox{$p$}{$p \to q$}{$q$} & $0$ & $1$ \\
        \hline
        $0$ & $1$ & $1$ \\
        \hline
        $1$ & $0$ & $1$ \\
        \hline
    \end{tabular}
    \caption{蕴涵的真值表}
\end{table}

蕴涵运算可以理解为，$p \to q$ 为真说明，若 $p$ 为真时，那么 $q$ 一定为真。

\subsubsection{互蕴涵/等价}

对于逻辑变量 $p, q$，互蕴涵运算记做 $p \leftrightarrow q$，其真值表如下：

\begin{table}[htbp]
    \centering
    \begin{tabular}{|c|c|c|}
        \hline
        \diagbox{$p$}{$p \leftrightarrow q$}{$q$} & $0$ & $1$ \\
        \hline
        $0$ & $1$ & $0$ \\
        \hline
        $1$ & $0$ & $1$ \\
        \hline
    \end{tabular}
    \caption{互蕴涵的真值表}
\end{table}

互蕴涵即 $(p \to q) \land (q \to p)$。

逻辑联结词的优先级为：非 > 合取 > 析取 > 异或 > 蕴涵 > 互蕴涵。
